% !TEX root = ../main.tex

\chapter{核心机制设计与实现}


\section{合约链上的 HTLC 实现}


在支持智能合约的链上，本文系统使用 HTLC 实现跨链资产锁定。合约中的 CrossChainSwap 结构记录用户侧交换状态，包括 hashLock、timeLock、sender、recipient、amount、completed、refunded、createdAt 和 isManagerLocked 等字段。ManagerLock 结构记录 Manager 侧响应锁定状态，包括 hashLock、timeLock、recipient、amount、completed、refunded 和 createdAt 等字段。


当 Alice 调用 initiateCrossChain 发起交换时，合约首先检查 hashLock 是否为空、recipient 是否有效、amount 是否大于零以及用户余额是否足够。随后合约生成 swapId，将 Alice 的资产转入合约地址锁定，并保存 CrossChainSwap 状态。该过程完成后，合约触发 CrossChainInitiated 事件，供 Notary 监听。


目标链 Manager 在收到 Notary 转发的请求后，调用 managerRespondToCrossChain 创建 ManagerLock。该函数要求调用者必须是 Manager，并要求 hashLock 未被重复使用。Manager 侧的锁定通常应设置比用户侧更短的 timeLock，这样可以减少目标链资产长时间占用的风险，也有利于失败时按合理顺序退款。


完成交易时，completeCrossChain 或 completeManagerLock 函数会要求调用方提交 secret。合约计算 keccak256(secret)，并与保存的 hashLock 比较。如果匹配，则将状态标记为 completed，并执行资产释放、销毁或结算逻辑。如果不匹配，交易会回滚。这一机制保证不知道 secret 的参与方无法领取资产。


退款时，refundCrossChain 和 refundManagerLock 依赖 timeLock 判断。只有超过规定时间且交易未完成、未退款时，退款逻辑才能执行。用户侧退款将锁定资产退回 Alice，Manager 侧退款则释放或回收 Manager 锁定的资产。通过完成和退款两条路径，合约形成了一个简明的跨链状态机。


\section{公证人协调流程}


公证人协调流程是连接源链和目标链的关键。由于不同链之间不能直接调用彼此合约，源链上的事件不会自动触发目标链状态变化，因此需要 Notary 监听源链事件并发起目标链交易。


在本文系统中，Notary 的典型流程包括：首先连接源链节点，订阅或轮询 CrossChainInitiated 事件；其次解析事件中的 swapId、hashLock、sender、recipient、amount 和 timeLock 等参数；然后根据 targetChainId 或系统配置确定目标链；最后调用目标链 Manager 或 TokenPool 合约创建对应锁定。Notary 还可以记录跨链状态，便于前端查询交易进度。


Notary 的设计需要注意信任边界。它可以延迟消息、漏转发消息或错误转发参数，但它不应拥有绕过 hashLock 直接释放资产的能力。目标链 ManagerLock 的完成仍然依赖正确 secret，源链退款仍然依赖 timeLock。因此，Notary 是流程协调者，不是最终安全性的唯一来源。安全性应由合约条件、签名协议和超时机制共同保证。


\section{无脚本链场景中的 adaptor signature}


对于 Monero、ZCash 屏蔽交易等无脚本链，链上不提供通用脚本或合约系统，HTLC 所依赖的哈希锁判断和条件分支无法在链上表达，因此无法直接迁移 HTLC。本文系统引入 Schnorr adaptor signature 作为这类链上的条件锁定方案。其核心过程是：参与方先构造一份与秘密值相关的预签名，该预签名本身不是有效签名，无法直接完成交易；当秘密值被揭示或补入后，预签名可以转化为正式签名，交易得以完成；同时，另一方可以根据预签名和正式签名之间的差异提取秘密。这样，原本需要链上脚本判断的"提交正确 secret 才能领取"条件，被等价地转化为"补全正确秘密才能生成有效签名"，仅凭签名能力即可完成。


这种设计与 HTLC 的目标一致，都是将两侧交易绑定到同一个秘密上。但二者的实现层不同：HTLC 在链上保存 hashLock，由合约或脚本验证 secret；adaptor signature 在链下签名协议中嵌入秘密条件，由签名补全关系保证交易完成与秘密公开之间的联系，链上只需验证一笔普通签名。正因如此，adaptor signature 不依赖目标链具备任何脚本能力，能够覆盖 HTLC 无法触及的无脚本链。


在本文原型中，scriptless 模块主要承担机制验证作用，以 BTC 作为无脚本链场景的替代实现。它模拟了 Alice、Bob 等参与方如何创建 adaptor signature，如何在链上观察签名完成，如何提取 secret，以及如何使用该 secret 完成另一侧交换。需要说明的是，选择 BTC 作为替代是出于其签名模型成熟、便于模拟的工程考虑，并不意味着 BTC 本身缺乏脚本能力；该模块验证的是 adaptor signature 在"仅有签名能力"假设下的正确性，这一正确性可直接迁移到 Monero、ZCash 屏蔽交易等真正的无脚本链。


\section{HTLC 与 adaptor signature 的统一关系}


HTLC 与 adaptor signature 并不是两套彼此割裂的跨链系统，而是统一跨链架构下针对不同链能力的两种实现方式。它们都服务于同一个目标：使跨链交易的完成依赖同一个秘密条件，并在失败时提供安全退出路径。两种机制的选择根植于目标链脚本表达能力的差异。


EVM 类链采用账户模型，具备较强的智能合约表达能力。合约可以维护全局状态映射（如 swapId → CrossChainSwap 结构），在链上直接保存 hashLock、timeLock、sender、recipient、amount 等字段，并通过合约函数显式执行锁定、条件领取和超时退款逻辑。整个跨链交换的状态机可以完整地在链上表达和验证，因此 HTLC 是合约链上自然且高效的选择。


无脚本链（如 Monero、ZCash 屏蔽交易）则处于另一个极端：它们以环签名、零知识证明等纯密码学方式构造交易，链上不提供通用脚本或合约系统，资产花费完全由签名授权决定。在这类链上，HTLC 所需的"提交正确 secret 才能领取"这一条件分支根本无法在链上表达，因此 HTLC 无法迁移。Schnorr adaptor signature 将条件锁定逻辑从链上脚本转移到链下签名协议中，交易的完成取决于签名是否被正确补全，而非脚本分支的执行路径。由于补全后的产物只是一笔普通签名，adaptor signature 不要求目标链具备任何脚本能力，从而把原子交换的适用范围从合约链扩展到了无脚本链。这一点是 adaptor signature 相对链上 HTLC 不可替代的价值：它解决的不是"是否更优"的问题，而是"在无脚本链上是否可行"的问题。


因此，本文系统的兼容性并不来自”所有链使用同一种机制”，而来自”所有链遵循同一角色流程，并根据自身脚本能力选择合适的底层机制”。在架构层，Alice、Manager、Notary、Bob 的流程保持一致；在机制层，合约链采用链上 HTLC，无脚本链采用 adaptor signature 以摆脱对脚本能力的依赖；在协调层，swapId、hashLock、secret、amount 和 recipient 构成统一上下文。
